% ==============================================================================
% Chapter 2 — Literature Review
% ==============================================================================

\chapter{Literature Review}
\label{ch:literature}

\todo[inline]{Write literature review.}

\section{European Emission Inventories}
\label{sec:lit_inventories}

Anthropogenic emission inventories are the primary input to air quality models.
At the European scale, the most widely used dataset is the \cams{} emission
inventory \citep{Kuenen2022}, which provides gridded anthropogenic emissions of
trace gases and particulate matter at a horizontal resolution of approximately
($\sim$\qtyproduct{6 x 6}{\km})
 for the pan-European domain.

\todo[inline]{Describe CAMS-REG structure: GNFR sectors, pollutants covered
(NOx, PM2.5, PM10, NMVOC, SO2, NH3), temporal resolution, base years available.}

\subsection{Road Transport Emissions}
\label{sec:lit_road}

\todo[inline]{Summarise improvements in the RI-URBANS version of CAMS-REG:
COPERT bottom-up model, cold start emissions, non-exhaust PM, OTM spatial
allocation and random-forest gapfilling \citep{Kuenen_ST15_2025}.}

\subsection{Ultrafine Particle Number Emissions}
\label{sec:lit_ufp}

\todo[inline]{Describe the TPN inventory approach: emission factors from
literature, vehicle type specificity, mass-based approach for stationary sources.}

\section{Urban Emission Downscaling Methods}
\label{sec:lit_downscaling}

Several approaches exist for spatial disaggregation of coarse regional
inventories to city-scale resolution. These range from simple proxy-based
methods using population density to hybrid approaches that combine multiple
source-specific spatial datasets \citep{Ramacher2021}.

\todo[inline]{Discuss MEGAPOLI, uEMEP \citep{Mu2022}, and UrbEm in detail.
Compare spatial proxies used, handling of line/point sources, mass consistency.}

\section{City-Scale Air Quality Modelling}
\label{sec:lit_aqm}

\todo[inline]{Introduce EPISODE-CityChem: Eulerian grid model + street network
module. Discuss input requirements (emission files: UECT format for points,
areas, lines). Cite relevant validation studies for Athens and Hamburg.}

\section{The Metronome Project and Study Cities}
\label{sec:lit_metronome}

\todo[inline]{Describe the HE Metronome project framework, the two Mission Cities
(Ioannina and Kozani), their urban air quality challenges, and the role of
emission modelling in the project.}
